\section{实现与可复现细节}
\label{app:protocol}
\subsection{图预处理与固定权重}
表~\ref{tab:datasets} 给出执行计划记录的数据维度。特征使用数据加载器的逐行归一化；官方训练与测试掩码分别定义支持池和查询集，验证集标签不用于拟合或选参。其他节点作为未标注图上下文可见。每个任务包含该图的全部类别，每类抽取指定数量的支持节点。
\begin{table}[!htbp]
\centering
\caption{\lang{Dataset dimensions and trainable prompt counts from the recorded protocol.}{记录协议中的数据维度及可学习提示参数量。}}
\label{tab:datasets}
\small
\setlength{\tabcolsep}{4pt}
\begin{tabular}{lrrrr}
\toprule
\lang{Dataset}{数据集} & \lang{Nodes}{节点} & \lang{Features}{特征} & \lang{Classes}{类别} & \lang{Prompt parameters}{提示参数} \\
\midrule
Cora & 2708 & 1433 & 7 & 896 \\
CiteSeer & 3327 & 3703 & 6 & 768 \\
PubMed & 19717 & 500 & 3 & 384 \\
\bottomrule
\end{tabular}
\end{table}


编码器由两层宽度为 128 的 GCNConv 构成，中间使用 LeakyReLU。第一层包含预训练仿射变换与图聚合，提示添加在该层输出、非线性激活之前；第二层卷积在所评估编码器中不再接额外激活。每个数据集的所有方法使用相同权重，结果记录保存权重、图及划分的哈希。预训练来源标识为 \texttt{Edgepred\_Gprompt}，预训练随机种子未记录。

\subsection{训练与预测细节}
正文描述实际生效的前向计算。Adam 使用 $\beta_1=0.9$、$\beta_2=0.999$、$\epsilon=10^{-8}$，学习率为 0.01，权重衰减为 0。提示正则为提示矩阵元素平方的均值乘以 $10^{-3}$，与优化器权重衰减不同。逐通道尺度使用分母为 $n-1$ 的样本标准差；欧氏归一化采用实现中的 $10^{-12}$ 范数下界。训练中没有生效的蒸馏目标或查询标签目标。

对于支持损失，令 $\bar z_i=\norm{Z_i}$，$m_c=\sum_{j\in S_c}\bar z_j$。当 $k>1$ 时，支持节点 $i$ 的训练原型为
\begin{equation}
p_c^{(-i)}=\norm{\frac{m_c-\mathbb I[y_i=c]\bar z_i}{k-\mathbb I[y_i=c]}},\qquad
\ell_{ic}=\bar z_i^\top p_c^{(-i)}/0.2.
\label{eq:loo}
\end{equation}
$k=1$ 时，以包含自身的类别原型 $\norm{m_c/k}$ 替代 $p_c^{(-i)}$。预测时定义
\begin{equation}
z_i^*=\norm{Z_i-\frac1n\sum_{j\in V}Z_j},\quad
p_c^*=\norm{\frac1k\sum_{i\in S_c}z_i^*},\quad
\widehat y_i=\arg\max_c (z_i^*)^\top p_c^*.
\label{eq:predict}
\end{equation}
将余弦 logits 除以 0.2 后取 softmax，得到五项指标使用的概率分数。训练时不使用预测阶段的输出中心化。读出是固定实现选择，不是本文的实验比较因素。

\paragraph{计算流程。}
给定图、固定预训练权重及支持标签：
\begin{enumerate}
\item 归一化邻接，计算式~\eqref{eq:attributes} 及第一层预激活表示。
\item 依式~\eqref{eq:anchor} 初始化支持原型，执行式~\eqref{eq:assignment}--\eqref{eq:refinement} 三轮。
\item 依式~\eqref{eq:routing} 计算固定路由；共享提示消融则将其替换为 $B_{ic}=1/C$。
\item 初始化 $P=0$ 并缓存隐藏通道尺度，固定编码器与路由。
\item 重复 100 步：运行提示编码器，构造支持 logits，依式~\eqref{eq:loss} 更新 $P$。
\item 依式~\eqref{eq:predict} 计算最终原型分数，封存预测后再使用查询标签评分。
\end{enumerate}

\subsection{基线拟合与随机种子}
支持抽样 seed 为 30093000 至 30093019，训练 seed 为 0、1、2。每次抽样在逐类支持池中生成有种子控制的排列，取前 $k$ 个节点；各方法接收相同支持身份。免训练的冻结预训练条件为对齐记录而保留三个 seed 槽位，不能将其视为三次独立拟合。

可学习比较方法使用支持交叉熵。全量微调与输入提示方法包含任务读出，GraphPrompt 适配实现则对提示后的表示进行原型预测。GPF 添加共享特征向量，GPF-plus 学习 20 个特征提示及输入相关组合。它们在 1-shot 下的固定参数为学习率 0.01、权重衰减 0.0001。3-shot 和 5-shot 的三个支持内折各留出每类一个节点，按留出支持交叉熵均值选择四种候选配置，再在完整支持集上拟合所选方法；不使用额外验证标签预算。该协议控制了标注信息与最终训练步数，但没有令总选参计算量或读出参数化完全一致。

\subsection{同一条件内的统计}
固定数据集、shot 数、预训练权重和方法，以 $a_{st}$ 表示支持 seed $s$、训练 seed $t$ 的某项指标，计算
\begin{equation}
\bar a_s=\frac13\sum_{t=1}^{3}a_{st},\quad
\widehat a=\frac1{20}\sum_{s=1}^{20}\bar a_s,\quad
\widehat\sigma=\sqrt{\frac1{19}\sum_{s=1}^{20}(\bar a_s-\widehat a)^2}.
\label{eq:stats}
\end{equation}
表中报告 $100\widehat a\pm100\widehat\sigma$，不是标准误。配对比较对 $d_s=\bar a_s^{\mathrm{PRPrompt}}-\bar a_s^{\mathrm{reference}}$ 按支持抽样块重采样，以 10,000 次重采样均值的百分位点构成 95\% 区间。原区间直接复用，不作修改。不跨数据集、shot 或方法计算指标均值与平均效应。

数据包保留 Accuracy、Macro-F1、Micro-F1、宏平均 one-vs-rest ROC-AUC 和宏平均 AP。单标签多分类下 Micro-F1 与 Accuracy 相同，此处为完整性一并列出。不同支持抽样共享固定查询集，区间反映该图条件下支持选择和训练的波动，并非独立图样本的不确定性。

\subsection{环境、来源与数值复核}
执行计划记录复用 ProG 环境中的 Python 3.11.16、PyTorch 2.11.0+cu130 与 CUDA 13.0，逐次结果记录 CUDA 设备及拟合时间。当前提取包未给出本报告范围独立核算的总 GPU 小时及已核实的加速器清单，本文不据此进行效率比较。

表格来源于封存的三数据集材料，包含主比较 3,240 个单元和额外共享提示 540 个单元，其中 \method 占 540 个单元。提取时核对原始指标文件校验和及配对图、权重、支持／查询身份。源数值审计显示，所选 3,780 个单元在保存表示上的读出全部精确复现，门控／前向误差通过阈值，新前向未产生查询类别变化。这些结论来自归档复核，不是本次重新运行实验；阈值通过也不等同于 CUDA 运算逐位一致。

代码与表格数据见随附仓库 \url{https://github.com/Mes-sala/PRPrompt}。已记录实验的源码版本为 \texttt{0328a2bb68f249a34c89557178b6233daf5ce20d}。论文包包含表格生成、静态检查脚本、哈希及文献来源记录。模型权重和优化器状态保留在实验服务器，不包含在本地论文包内。本文整理既有实验，没有为撰稿新增训练或 bootstrap 分析。

\section{其余逐条件指标}
\label{app:more}
以下表格按相同 seed 统计口径补充主比较的其余指标，每一行仍对应独立的数据集／shot 条件。
\begin{table}[htbp]
\centering
\caption{\lang{Micro-F1 (\%, higher is better). Main comparison. Within each dataset/shot/method, training seeds are averaged within each support draw; mean and sample standard deviation are then computed over 20 support draws. Bold marks the largest mean in that row. $\dagger$: common-protocol adaptation.}{Micro-F1（\%，越高越好）。主比较。固定数据集、shot 与方法，在每个支持抽样内先平均训练 seed，再对 20 个支持抽样计算均值与样本标准差。粗体为该行最大均值。 $\dagger$：统一协议适配实现。}}
\label{tab:main-micro_f1}
\small
\setlength{\tabcolsep}{4pt}
\resizebox{\linewidth}{!}{%
\begin{tabular}{lrrrrrr}
\toprule
\lang{Dataset / shot}{数据集 / shot} & \frozenname & \ftname & GPF & GPF-plus & \gpromptname & \method \\
\midrule
Cora / 1 & \meanstd{23.17}{5.48} & \meanstd{45.41}{10.18} & \meanstd{15.20}{3.24} & \meanstd{30.80}{6.64} & \meanstd{21.88}{4.62} & \bestmeanstd{54.59}{9.18} \\
Cora / 3 & \meanstd{29.42}{3.47} & \meanstd{61.54}{6.78} & \meanstd{24.32}{4.66} & \meanstd{40.21}{4.80} & \meanstd{27.75}{3.64} & \bestmeanstd{70.76}{2.76} \\
Cora / 5 & \meanstd{33.56}{5.03} & \meanstd{67.50}{5.06} & \meanstd{28.17}{5.54} & \meanstd{53.14}{5.55} & \meanstd{31.36}{3.61} & \bestmeanstd{73.02}{1.95} \\
CiteSeer / 1 & \meanstd{21.14}{2.87} & \meanstd{35.76}{6.80} & \meanstd{17.10}{2.64} & \meanstd{21.40}{4.60} & \meanstd{19.53}{3.34} & \bestmeanstd{47.12}{6.31} \\
CiteSeer / 3 & \meanstd{25.29}{2.16} & \meanstd{42.22}{5.88} & \meanstd{20.36}{2.02} & \meanstd{28.92}{2.68} & \meanstd{24.83}{2.77} & \bestmeanstd{61.20}{4.98} \\
CiteSeer / 5 & \meanstd{27.38}{2.84} & \meanstd{47.89}{4.85} & \meanstd{21.40}{2.57} & \meanstd{31.82}{2.60} & \meanstd{27.66}{2.49} & \bestmeanstd{66.17}{2.16} \\
PubMed / 1 & \meanstd{55.44}{11.42} & \meanstd{48.77}{8.54} & \meanstd{36.29}{5.79} & \meanstd{41.12}{6.83} & \meanstd{50.24}{12.95} & \bestmeanstd{58.70}{9.78} \\
PubMed / 3 & \meanstd{65.55}{8.11} & \meanstd{59.75}{7.97} & \meanstd{51.28}{6.93} & \meanstd{57.46}{4.96} & \meanstd{62.47}{5.54} & \bestmeanstd{68.66}{4.58} \\
PubMed / 5 & \meanstd{68.37}{3.41} & \meanstd{63.96}{5.18} & \meanstd{55.83}{4.93} & \meanstd{61.37}{5.32} & \meanstd{66.94}{4.09} & \bestmeanstd{70.98}{2.46} \\
\bottomrule
\end{tabular}
}
\end{table}

\begin{table}[!htbp]
\centering
\caption{\lang{Macro ROC-AUC (\%, higher is better). Main comparison; seed aggregation follows Appendix~\ref{app:protocol}.}{Macro ROC-AUC（\%，越高越好）。主比较；seed 统计沿用附录~\ref{app:protocol}。}}
\label{tab:main-roc_auc_macro}
\small
\setlength{\tabcolsep}{4pt}
\resizebox{\linewidth}{!}{%
\begin{tabular}{lrrrrrr}
\toprule
\lang{Dataset / shot}{数据集 / shot} & \frozenname & \ftname & GPF & GPF-plus & \gpromptname & \method \\
\midrule
Cora / 1 & \meanstd{62.95}{4.93} & \meanstd{81.83}{5.73} & \meanstd{53.20}{2.50} & \meanstd{74.73}{4.00} & \meanstd{61.93}{4.85} & \bestmeanstd{83.06}{5.42} \\
Cora / 3 & \meanstd{68.95}{2.99} & \meanstd{89.94}{2.35} & \meanstd{64.61}{4.25} & \meanstd{78.13}{2.49} & \meanstd{68.05}{3.01} & \bestmeanstd{90.84}{1.97} \\
Cora / 5 & \meanstd{72.36}{2.85} & \bestmeanstd{92.58}{1.26} & \meanstd{68.65}{4.34} & \meanstd{85.66}{2.71} & \meanstd{71.01}{2.66} & \meanstd{92.47}{1.11} \\
CiteSeer / 1 & \meanstd{56.11}{3.24} & \meanstd{70.97}{4.63} & \meanstd{53.08}{1.84} & \meanstd{64.19}{4.07} & \meanstd{54.39}{4.21} & \bestmeanstd{74.71}{5.81} \\
CiteSeer / 3 & \meanstd{60.17}{2.06} & \meanstd{76.04}{3.00} & \meanstd{55.85}{1.91} & \meanstd{65.37}{3.23} & \meanstd{59.99}{2.60} & \bestmeanstd{82.82}{3.08} \\
CiteSeer / 5 & \meanstd{62.69}{2.35} & \meanstd{79.54}{2.22} & \meanstd{57.93}{1.77} & \meanstd{68.91}{2.38} & \meanstd{62.67}{2.33} & \bestmeanstd{85.86}{1.52} \\
PubMed / 1 & \meanstd{73.87}{11.03} & \meanstd{70.73}{9.76} & \meanstd{60.13}{8.31} & \meanstd{68.63}{6.61} & \meanstd{70.43}{11.06} & \bestmeanstd{76.48}{10.43} \\
PubMed / 3 & \meanstd{82.17}{6.59} & \meanstd{81.00}{3.74} & \meanstd{72.19}{4.91} & \meanstd{76.82}{4.24} & \meanstd{81.13}{4.42} & \bestmeanstd{86.29}{3.10} \\
PubMed / 5 & \meanstd{84.83}{2.59} & \meanstd{84.31}{2.45} & \meanstd{74.64}{4.10} & \meanstd{79.98}{2.95} & \meanstd{83.57}{2.81} & \bestmeanstd{87.03}{1.87} \\
\bottomrule
\end{tabular}
}
\end{table}

\begin{table}[!htbp]
\centering
\caption{\lang{Macro AP (\%, higher is better). Main comparison; seed aggregation follows Appendix~\ref{app:protocol}.}{Macro AP（\%，越高越好）。主比较；seed 统计沿用附录~\ref{app:protocol}。}}
\label{tab:main-average_precision_macro}
\small
\setlength{\tabcolsep}{4pt}
\resizebox{\linewidth}{!}{%
\begin{tabular}{lrrrrrr}
\toprule
\lang{Dataset / shot}{数据集 / shot} & \frozenname & \ftname & GPF & GPF-plus & \gpromptname & \method \\
\midrule
Cora / 1 & \meanstd{23.64}{4.26} & \meanstd{54.01}{10.10} & \meanstd{16.64}{1.47} & \meanstd{40.37}{5.88} & \meanstd{23.15}{3.89} & \bestmeanstd{59.04}{9.07} \\
Cora / 3 & \meanstd{29.20}{3.71} & \meanstd{70.20}{4.98} & \meanstd{24.89}{3.11} & \meanstd{44.43}{4.09} & \meanstd{27.06}{2.89} & \bestmeanstd{74.33}{3.38} \\
Cora / 5 & \meanstd{33.16}{4.26} & \meanstd{75.99}{3.16} & \meanstd{28.68}{3.71} & \meanstd{58.84}{6.14} & \meanstd{29.92}{3.11} & \bestmeanstd{77.64}{1.87} \\
CiteSeer / 1 & \meanstd{21.53}{2.32} & \meanstd{38.74}{5.84} & \meanstd{19.66}{1.43} & \meanstd{29.38}{4.20} & \meanstd{20.61}{2.27} & \bestmeanstd{45.99}{6.79} \\
CiteSeer / 3 & \meanstd{24.17}{1.59} & \meanstd{45.56}{4.79} & \meanstd{21.40}{1.15} & \meanstd{30.59}{3.23} & \meanstd{23.97}{1.89} & \bestmeanstd{59.15}{4.80} \\
CiteSeer / 5 & \meanstd{26.46}{2.07} & \meanstd{51.03}{4.40} & \meanstd{22.88}{1.45} & \meanstd{34.40}{2.70} & \meanstd{26.04}{1.65} & \bestmeanstd{64.02}{2.66} \\
PubMed / 1 & \meanstd{57.52}{10.93} & \meanstd{55.36}{10.55} & \meanstd{43.93}{6.89} & \meanstd{52.61}{8.08} & \meanstd{52.39}{9.87} & \bestmeanstd{61.58}{11.39} \\
PubMed / 3 & \meanstd{67.27}{8.59} & \meanstd{67.52}{6.82} & \meanstd{55.29}{5.93} & \meanstd{60.88}{5.58} & \meanstd{65.03}{5.79} & \bestmeanstd{74.22}{5.19} \\
PubMed / 5 & \meanstd{70.34}{4.13} & \meanstd{72.33}{5.00} & \meanstd{58.42}{5.12} & \meanstd{65.42}{4.92} & \meanstd{68.42}{3.67} & \bestmeanstd{75.35}{3.34} \\
\bottomrule
\end{tabular}
}
\end{table}

\clearpage
\section{配对差值与不确定性}
\label{app:paired}
各单元格报告 \method 减去指定对照的差值，单位为百分点，并附未经多重比较校正的 95\% 配对支持 bootstrap 区间。所有数值保留原配对记录，表格仅组织不同条件，不对其求平均。
\begin{table}[!htbp]
\centering
\caption{\lang{\method{} minus \frozenname: paired differences and 95\% intervals (percentage points).}{\method{} 相对 \frozenname 的配对差值与 95\% 区间（百分点）。}}
\label{tab:paired-FrozenPrototype}
\small
\setlength{\tabcolsep}{4pt}
\begin{tabular}{lrr}
\toprule
\lang{Dataset / shot}{数据集 / shot} & Accuracy & Macro-F1 \\
\midrule
Cora / 1 & $+31.42\;[+28.20,\,+34.59]$ & $+30.84\;[+27.77,\,+33.84]$ \\
Cora / 3 & $+41.34\;[+39.63,\,+43.04]$ & $+40.32\;[+38.64,\,+42.08]$ \\
Cora / 5 & $+39.46\;[+37.23,\,+41.61]$ & $+38.97\;[+36.88,\,+40.99]$ \\
CiteSeer / 1 & $+25.98\;[+23.72,\,+28.21]$ & $+23.87\;[+21.80,\,+25.77]$ \\
CiteSeer / 3 & $+35.91\;[+33.05,\,+38.39]$ & $+33.87\;[+31.34,\,+36.14]$ \\
CiteSeer / 5 & $+38.78\;[+37.37,\,+40.18]$ & $+36.58\;[+35.22,\,+37.95]$ \\
PubMed / 1 & $+3.26\;[+0.40,\,+6.22]$ & $+3.93\;[+0.86,\,+7.00]$ \\
PubMed / 3 & $+3.11\;[+0.38,\,+6.27]$ & $+5.02\;[+2.27,\,+8.44]$ \\
PubMed / 5 & $+2.61\;[+1.13,\,+4.22]$ & $+4.12\;[+2.53,\,+5.91]$ \\
\bottomrule
\end{tabular}
\end{table}

\begin{table}[!htbp]
\centering
\caption{\lang{\method{} minus \ftname: paired differences and 95\% intervals (percentage points).}{\method{} 相对 \ftname 的配对差值与 95\% 区间（百分点）。}}
\label{tab:paired-FullFineTune}
\small
\setlength{\tabcolsep}{4pt}
\begin{tabular}{lrr}
\toprule
\lang{Dataset / shot}{数据集 / shot} & Accuracy & Macro-F1 \\
\midrule
Cora / 1 & $+9.17\;[+6.58,\,+11.87]$ & $+8.97\;[+6.40,\,+11.77]$ \\
Cora / 3 & $+9.22\;[+6.75,\,+11.79]$ & $+8.44\;[+6.12,\,+10.90]$ \\
Cora / 5 & $+5.52\;[+3.73,\,+7.50]$ & $+5.05\;[+3.34,\,+6.86]$ \\
CiteSeer / 1 & $+11.36\;[+9.57,\,+13.18]$ & $+13.27\;[+11.17,\,+15.31]$ \\
CiteSeer / 3 & $+18.98\;[+16.78,\,+20.94]$ & $+19.76\;[+17.45,\,+22.00]$ \\
CiteSeer / 5 & $+18.28\;[+16.22,\,+20.13]$ & $+18.42\;[+16.07,\,+20.60]$ \\
PubMed / 1 & $+9.93\;[+6.77,\,+13.14]$ & $+12.86\;[+9.39,\,+16.37]$ \\
PubMed / 3 & $+8.90\;[+5.89,\,+12.26]$ & $+10.57\;[+7.08,\,+14.41]$ \\
PubMed / 5 & $+7.02\;[+4.96,\,+9.14]$ & $+7.86\;[+5.54,\,+10.28]$ \\
\bottomrule
\end{tabular}
\end{table}

\clearpage
\begin{table}[!htbp]
\centering
\caption{\lang{\method{} minus GPF: paired differences and 95\% intervals (percentage points).}{\method{} 相对 GPF 的配对差值与 95\% 区间（百分点）。}}
\label{tab:paired-GPF}
\small
\setlength{\tabcolsep}{4pt}
\begin{tabular}{lrr}
\toprule
\lang{Dataset / shot}{数据集 / shot} & Accuracy & Macro-F1 \\
\midrule
Cora / 1 & $+39.38\;[+35.60,\,+42.98]$ & $+46.27\;[+42.36,\,+49.82]$ \\
Cora / 3 & $+46.44\;[+44.21,\,+48.74]$ & $+48.25\;[+45.22,\,+51.53]$ \\
Cora / 5 & $+44.86\;[+42.52,\,+47.57]$ & $+46.17\;[+43.71,\,+49.32]$ \\
CiteSeer / 1 & $+30.02\;[+27.04,\,+33.00]$ & $+37.48\;[+34.93,\,+39.92]$ \\
CiteSeer / 3 & $+40.83\;[+38.11,\,+43.06]$ & $+44.38\;[+40.82,\,+47.90]$ \\
CiteSeer / 5 & $+44.77\;[+43.32,\,+46.16]$ & $+47.77\;[+44.95,\,+50.72]$ \\
PubMed / 1 & $+22.40\;[+18.82,\,+26.24]$ & $+31.76\;[+28.23,\,+35.27]$ \\
PubMed / 3 & $+17.37\;[+14.16,\,+21.02]$ & $+21.33\;[+17.31,\,+25.89]$ \\
PubMed / 5 & $+15.15\;[+13.30,\,+17.01]$ & $+17.97\;[+15.46,\,+20.67]$ \\
\bottomrule
\end{tabular}
\end{table}

\begin{table}[!htbp]
\centering
\caption{\lang{\method{} minus GPF-plus: paired differences and 95\% intervals (percentage points).}{\method{} 相对 GPF-plus 的配对差值与 95\% 区间（百分点）。}}
\label{tab:paired-GPFPlus}
\small
\setlength{\tabcolsep}{4pt}
\begin{tabular}{lrr}
\toprule
\lang{Dataset / shot}{数据集 / shot} & Accuracy & Macro-F1 \\
\midrule
Cora / 1 & $+23.78\;[+19.96,\,+27.47]$ & $+26.26\;[+22.23,\,+30.20]$ \\
Cora / 3 & $+30.55\;[+28.29,\,+32.52]$ & $+30.07\;[+28.04,\,+31.97]$ \\
Cora / 5 & $+19.88\;[+17.62,\,+22.16]$ & $+20.22\;[+18.02,\,+22.39]$ \\
CiteSeer / 1 & $+25.71\;[+22.89,\,+28.45]$ & $+31.43\;[+28.61,\,+34.12]$ \\
CiteSeer / 3 & $+32.27\;[+30.45,\,+34.03]$ & $+32.03\;[+30.10,\,+33.81]$ \\
CiteSeer / 5 & $+34.35\;[+33.05,\,+35.65]$ & $+33.57\;[+32.22,\,+34.87]$ \\
PubMed / 1 & $+17.57\;[+13.28,\,+21.86]$ & $+22.48\;[+18.10,\,+26.98]$ \\
PubMed / 3 & $+11.20\;[+9.47,\,+13.11]$ & $+13.74\;[+12.07,\,+15.56]$ \\
PubMed / 5 & $+9.61\;[+7.84,\,+11.49]$ & $+11.61\;[+9.80,\,+13.53]$ \\
\bottomrule
\end{tabular}
\end{table}

\clearpage
\begin{table}[!htbp]
\centering
\caption{\lang{\method{} minus \gpromptname: paired differences and 95\% intervals (percentage points).}{\method{} 相对 \gpromptname 的配对差值与 95\% 区间（百分点）。}}
\label{tab:paired-GraphPromptDiag}
\small
\setlength{\tabcolsep}{4pt}
\begin{tabular}{lrr}
\toprule
\lang{Dataset / shot}{数据集 / shot} & Accuracy & Macro-F1 \\
\midrule
Cora / 1 & $+32.70\;[+29.19,\,+36.29]$ & $+32.32\;[+28.93,\,+35.68]$ \\
Cora / 3 & $+43.01\;[+41.38,\,+44.73]$ & $+42.41\;[+40.61,\,+44.20]$ \\
Cora / 5 & $+41.66\;[+40.03,\,+43.20]$ & $+41.22\;[+39.64,\,+42.75]$ \\
CiteSeer / 1 & $+27.59\;[+24.97,\,+30.09]$ & $+26.26\;[+23.90,\,+28.62]$ \\
CiteSeer / 3 & $+36.37\;[+33.26,\,+38.82]$ & $+34.48\;[+31.80,\,+36.73]$ \\
CiteSeer / 5 & $+38.51\;[+37.14,\,+39.90]$ & $+36.31\;[+34.99,\,+37.65]$ \\
PubMed / 1 & $+8.45\;[+3.60,\,+13.47]$ & $+12.02\;[+7.33,\,+16.52]$ \\
PubMed / 3 & $+6.19\;[+4.04,\,+8.30]$ & $+8.27\;[+6.21,\,+10.36]$ \\
PubMed / 5 & $+4.03\;[+1.94,\,+6.36]$ & $+5.60\;[+3.49,\,+7.88]$ \\
\bottomrule
\end{tabular}
\end{table}

\begin{table}[!htbp]
\centering
\caption{\lang{\method{} minus \sharedname: paired differences and 95\% intervals (percentage points).}{\method{} 相对 \sharedname 的配对差值与 95\% 区间（百分点）。}}
\label{tab:paired-NoRouting}
\small
\setlength{\tabcolsep}{4pt}
\begin{tabular}{lrr}
\toprule
\lang{Dataset / shot}{数据集 / shot} & Accuracy & Macro-F1 \\
\midrule
Cora / 1 & $+31.96\;[+28.59,\,+35.21]$ & $+31.70\;[+28.29,\,+35.00]$ \\
Cora / 3 & $+39.06\;[+36.49,\,+41.44]$ & $+39.13\;[+36.60,\,+41.48]$ \\
Cora / 5 & $+36.94\;[+34.87,\,+39.06]$ & $+37.15\;[+35.09,\,+39.23]$ \\
CiteSeer / 1 & $+26.40\;[+24.07,\,+28.73]$ & $+24.69\;[+22.70,\,+26.65]$ \\
CiteSeer / 3 & $+37.18\;[+34.90,\,+39.33]$ & $+35.52\;[+33.31,\,+37.55]$ \\
CiteSeer / 5 & $+40.45\;[+39.06,\,+41.84]$ & $+38.48\;[+37.22,\,+39.72]$ \\
PubMed / 1 & $+4.85\;[+1.26,\,+8.26]$ & $+5.44\;[+1.84,\,+9.07]$ \\
PubMed / 3 & $+5.69\;[+3.25,\,+8.27]$ & $+7.27\;[+4.64,\,+10.29]$ \\
PubMed / 5 & $+3.02\;[+1.39,\,+4.51]$ & $+4.10\;[+2.56,\,+5.58]$ \\
\bottomrule
\end{tabular}
\end{table}

